% ============================================================
%  Predicción óptima de marcadores en el Mundial 2026:
%  integrando mercados de predicción con modelos bayesianos
%  adaptativos
%
%  Revista de Matemática: Teoría y Aplicaciones (UCR/CIMPA)
%  Autor: Jorge A. Alfaro-Murillo
%  Compilar: pdflatex -synctex=1 main && bibtex main && pdflatex -synctex=1 main && pdflatex -synctex=1 main
% ============================================================

\documentclass[12pt]{article}

\usepackage[utf8]{inputenc}
\usepackage[T1]{fontenc}
\usepackage[spanish]{babel}
\usepackage{natbib}
\usepackage{amsmath, amssymb, amsthm}
\usepackage{booktabs}
\usepackage{hyperref}
\usepackage{graphicx}
\usepackage{xcolor}
\usepackage{geometry}
\geometry{margin=2.5cm}

% Para tabla de marcadores
\usepackage{array}
\usepackage{multirow}

% URLs en referencias
\usepackage{url}

% ── Entornos matemáticos ───────────────────────────────────────
\theoremstyle{remark}
\newtheorem{observacion}{Observación}

% ── Comandos de autor ──────────────────────────────────────────
\newcommand{\TODO}[1]{\textbf{\color{red}[TODO: #1]}}
\newcommand{\E}{\mathbb{E}}
\newcommand{\Prob}{\mathbb{P}}

% ── Información del artículo ───────────────────────────────────
\title{%
  \textbf{Predicción óptima de marcadores en el Mundial 2026:
  integrando mercados de predicción con modelos bayesianos
  adaptativos}\\[6pt]
  \large\textit{Optimal scoreline prediction for the 2026 World Cup:
  integrating prediction markets with adaptive Bayesian models}
}

\author{%
  Jorge A.\ Alfaro-Murillo\\
  \small Centro de Investigación en Matemática Pura y Aplicada (CIMPA)\\
  \small Universidad de Costa Rica\\
  \small San José, Costa Rica\\
  \small \href{mailto:jorge.alfaro@ucr.ac.cr}{jorge.alfaro@ucr.ac.cr}
}

\date{Versión: \today}

\begin{document}

\maketitle

% ── MSC 2020 ──────────────────────────────────────────────────
\noindent\textbf{MSC 2020:} 62F15, 62P25, 60G25, 91B82.

\noindent\textbf{Palabras clave:} mercados de predicción,
inferencia bayesiana, distribución de Poisson, predicción de
marcadores, fútbol, Mundial 2026, quiniela, optimización.

\noindent\textbf{Keywords:} prediction markets, Bayesian inference,
Poisson distribution, scoreline prediction, soccer, World Cup 2026,
betting pool, optimization.

\bigskip\hrule\bigskip

% ─────────────────────────────────────────────────────────────
\begin{abstract}
En una quiniela de fútbol, acertar el marcador exacto puede valer
hasta cinco veces más que simplemente predecir el ganador.
Esta asimetría en la regla de puntuación convierte la predicción
de marcadores en un problema de optimización matemático no trivial:
el marcador más probable no es necesariamente el más rentable.

Este artículo presenta un modelo estadístico para el Mundial de
Fútbol 2026 que maximiza los puntos esperados en una quiniela de
marcadores exactos mediante la integración de dos fuentes de
información complementarias.
La primera fuente son los \emph{mercados de predicción}:
Kalshi, una plataforma estadounidense regulada donde miles de
participantes negocian contratos sobre cada resultado del torneo.
Los precios de mercado agregan información dispersa de manera
eficiente, ofreciendo estimaciones de probabilidad que superan
en calibración a los pronósticos individuales de expertos.
La segunda fuente son tres ediciones del Mundial (2014, 2018 y 2022),
cuya distribución histórica de marcadores captura la estructura
estadística del juego a nivel de selecciones nacionales.

Ambas fuentes se integran mediante un modelo bayesiano adaptativo.
Los parámetros de decaimiento temporal~$\gamma$ ---que mide cuánto
pesa el historial reciente sobre el antiguo--- y de adaptación al
torneo en curso~$\delta$ ---que mide cuánto peso adicional reciben
los resultados del propio WC 2026--- se estiman por validación cruzada:
$\gamma \approx 0.84$ pondera más el historial reciente,
mientras que~$\delta$ crece conforme se acumulan resultados del
propio Mundial 2026, reflejando las características propias del
presente mundial como el número de equipos, el formato de grupos
y el nivel general de los seleccionados participantes.
Sobre la distribución resultante de marcadores se aplica un argumento
de optimización explícito para encontrar la predicción que maximiza
los puntos esperados bajo la regla de puntuación de la quiniela.

\TODO{Completar con resultados numéricos al concluir la fase de grupos.}

El código fuente del modelo y la plataforma web de seguimiento
están disponibles públicamente en GitHub.
\end{abstract}

\bigskip

% ── Abstract en inglés ────────────────────────────────────────
\noindent\textbf{Abstract.}
In a soccer quiniela, predicting the exact scoreline can be worth
up to five times more than simply predicting the winner.
This scoring-rule asymmetry turns scoreline prediction into a
non-trivial mathematical optimization problem: the most likely
scoreline is not necessarily the most valuable one.

This paper presents a statistical model for the 2026 FIFA World Cup
that maximizes expected points in an exact-scoreline betting pool
by integrating two complementary information sources.
The first is \emph{prediction markets}: Kalshi, a regulated
U.S.\ platform where thousands of participants trade contracts on each
tournament match.
Market prices efficiently aggregate dispersed information, providing
probability estimates that outperform individual expert forecasts
in calibration.
The second source is three editions of the World Cup (2014, 2018,
2022), whose historical scoreline distributions capture the
statistical structure of international tournament soccer.

Both sources are integrated via an adaptive Bayesian model.
Temporal decay $\gamma$ ---measuring how much recent tournaments
outweigh older ones--- and current-tournament adaptation $\delta$
---measuring the additional weight given to WC 2026 results---
are estimated by cross-validation: $\gamma \approx 0.84$
upweights recent history, while $\delta$ grows as 2026 results
accumulate, reflecting the specific characteristics of the current
tournament such as the number of teams, group format, and overall
level of the participating squads.
An explicit optimization argument then identifies the scoreline
prediction that maximizes expected points under the quiniela scoring
rule.

\TODO{Insert numerical results upon conclusion of the group stage.}

Model source code and live tracking website are publicly available
on GitHub.

\bigskip\hrule

% ─────────────────────────────────────────────────────────────
\section{Introducción}
\label{sec:intro}

La quiniela de fútbol es un fenómeno cultural ampliamente arraigado
en América Latina y España \citep{aguero2024, lopez2022, bravo2026,
lanacion2026arg}.
En su variante de marcador exacto, los participantes pronostican el
resultado preciso de cada partido de un torneo, acumulando puntos
según la cercanía de su predicción al resultado real.
La regla de puntuación habitual premia de manera desigual los distintos
tipos de acierto, por ejemplo: acertar el marcador exacto vale cinco puntos,
mientras que acertar únicamente al ganador vale solo dos.
Esta asimetría introduce un problema de optimización que va más allá de
simplemente predecir el resultado más probable: el marcador modal
puede diferir del marcador que maximiza los puntos esperados.

Formalicemos el problema.
Sea $P(h, k)$ la distribución de probabilidad sobre los marcadores
posibles de un partido, donde $h$ es el número de goles de un equipo
y $k$ el del otro.
Dado un sistema de puntuación $f(a, b, h, k)$ que asigna puntos
cuando se predice $(a, b)$ y el marcador real es $(h, k)$, la
predicción óptima es
\[
  (\hat{a}, \hat{b})
  = \arg\max_{(a,b)} \E[\text{pts} \mid \text{predicción}(a,b)]
  = \arg\max_{(a,b)} \sum_{h,k} f(a, b, h, k) \cdot P(h, k).
\]
El desafío central es estimar $P(h, k)$ con precisión para cada
partido del torneo.

\subsection*{Fuentes de información}

Dos fuentes de información son naturalmente complementarias para
este propósito:

\textbf{Historial de Mundiales.}
Las ediciones recientes del Mundial de Fútbol (2014, 2018 y 2022)
proveen 192 partidos cuyos marcadores capturan la distribución
estadística del juego entre selecciones nacionales de primer nivel.
Esta distribución refleja la estructura del fútbol de alto rendimiento:
la predominancia de marcadores bajos, la asimetría entre goles en
tiempo regular y tiempo extra, y la diferencia entre la fase de grupos
y la fase eliminatoria.

\textbf{Mercados de predicción.}
Los mercados de predicción son plataformas donde los participantes
negocian contratos cuyo valor depende del resultado de eventos futuros.
En el equilibrio, los precios de mercado convergen hacia las
probabilidades colectivas del resultado, agregando la información
dispersa de todos los participantes de manera análoga a cómo los
mercados financieros agregan expectativas sobre el valor futuro de
los activos \citep{arrow2008}.
La evidencia empírica muestra que estos mercados son competitivos con
las mejores herramientas de pronóstico disponibles \citep{spann2009}.

Para el Mundial 2026 utilizamos \textbf{Kalshi}
(\href{https://kalshi.com}{kalshi.com}), una plataforma de mercados
de predicción regulada en Estados Unidos.
Para cada partido del torneo, Kalshi ofrece tres tipos de mercados:
resultado del partido (victoria local, empate, victoria visitante),
total de goles y ventaja de goles, cada uno con alta liquidez y
precios actualizados en tiempo real.

\subsection*{Contribución}

Este artículo hace tres contribuciones.
Primero, integra precios de mercados de predicción de Kalshi con un
modelo bayesiano de distribución de marcadores para obtener
estimaciones $P(h,k)$ más precisas que cualquiera de las dos fuentes
por separado.
Segundo, aplica un argumento de optimización explícito bajo la regla
de puntuación de quiniela para identificar la predicción de marcador
que maximiza los puntos esperados.
Tercero, implementa el modelo completo en una plataforma web pública
con actualizaciones automáticas, y documenta el papel que jugó
Claude Code —una herramienta de programación asistida por inteligencia
artificial— en el desarrollo del sistema.

Hasta donde sabemos, esta es la primera publicación académica que
combina explícitamente precios de mercados de predicción con modelos
bayesianos informados por datos históricos de torneos anteriores y
actualizados secuencialmente conforme avanza el torneo en curso,
con el objetivo de optimizar pronósticos en concursos de predicción
deportiva con reglas de puntuación asimétricas.

% ─────────────────────────────────────────────────────────────
\section{Mercados de predicción}
\label{sec:mercados}

\subsection{Teoría y eficiencia}

Un mercado de predicción es un mecanismo de agregación de
información en el que los participantes negocian contratos
contingentes sobre eventos futuros.
Un contrato binario típico paga \$1 si el evento ocurre y \$0
en caso contrario; su precio de mercado es, en equilibrio, la
estimación colectiva de la probabilidad del evento \citep{surowiecki2004}.

La eficiencia de los mercados de predicción descansa sobre un
argumento de arbitraje: si el precio de mercado difiere
sustancialmente de la probabilidad verdadera, existe una oportunidad
de ganancia que agentes racionales explotarán hasta eliminarla.
Este mecanismo es el mismo que da eficiencia a los mercados
financieros clásicos, con la ventaja adicional de que los contratos
tienen un horizonte de tiempo definido y una resolución objetiva.

Empíricamente, una comparación para 678 partidos de la Bundesliga
alemana muestra que los mercados de predicción son al menos tan
precisos como las cuotas de bookmakers, y ambos superan claramente
a los pronosticadores individuales \citep{spann2009}.
Resultados similares se obtienen para el fútbol inglés medidos
con el score de Brier \citep{forrest2005}.

\subsection{Kalshi: arquitectura y mercados disponibles}

Kalshi es una plataforma de mercados de predicción regulada por la
Commodity Futures Trading Commission (CFTC) de Estados Unidos,
fundada en 2018 y operativa desde 2021.
Para el Mundial 2026, Kalshi ofrece tres familias de mercados
por partido:

\begin{enumerate}
  \item \textbf{Resultado del partido} (\texttt{KXWCGAME-}):
    submercados de victoria local, empate y victoria visitante.
  \item \textbf{Total de goles} (\texttt{KXWCTOTAL-}):
    submercados del tipo ``más de $n - 0.5$ goles'', para
    $n = 1, 2, \ldots, 6$.
  \item \textbf{Ventaja de goles} (\texttt{KXWCSPREAD-}):
    submercados del tipo ``el equipo $X$ gana por más de $m - 0.5$
    goles'', para distintos valores de $m$.
\end{enumerate}

Los precios se obtienen como el punto medio entre la mejor oferta de
compra y la mejor oferta de venta (midprice) cuando el spread es menor
de \$0.30; de lo contrario se descarta el mercado por falta de liquidez.

\subsection{Sesgo favorito-longshot}

Un hallazgo consistente en la literatura de mercados deportivos es el
\emph{sesgo favorito-longshot}: los contratos sobre resultados
probables (favoritos) están subvalorados, mientras que los contratos
sobre resultados improbables (longshots) están sobrevalorados
\citep{whelan2024}.
Un análisis de más de 300{,}000 contratos de Kalshi documenta
este sesgo: los contratos baratos ganan menos de lo implícito en
su precio, y los contratos caros generan pequeños retornos
positivos \citep{burgi2026}.
El modelo incorpora una corrección explícita mediante una
transformación de potencia que se describe en la Sección~\ref{sec:kalshi}.

\subsection{Plataformas alternativas}

Otras plataformas relevantes de mercados de predicción son:

\textbf{Manifold Markets} (\href{https://manifold.markets}{manifold.markets}):
plataforma de acceso global con dinero de juego (no real), con más de
un millón de mercados sobre eventos de todo tipo.
La ausencia de dinero real reduce los incentivos de corrección, pero
también elimina barreras regulatorias y permite mercados sobre
preguntas de nicho.

\textbf{Metaculus} (\href{https://metaculus.com}{metaculus.com}):
plataforma de pronóstico basada en torneos y reputación.
Los datos del proyecto Calibration City \citep{manifund2024}
sugieren que Metaculus obtiene scores de Brier de 0.084--0.111,
comparables con los mejores forecasters individuales.

\textbf{Polymarket}: plataforma de dinero real basada en blockchain,
no disponible para usuarios en Estados Unidos.
En las elecciones presidenciales de EE.UU.\ de 2024, Polymarket
superó a las encuestas tradicionales en precisión \citep{cutting2025}.

Para el Mundial 2026, Kalshi es la opción más adecuada por su
liquidez en mercados deportivos, su regulación en EE.UU.\ y la
disponibilidad de una API pública sin autenticación.

% ─────────────────────────────────────────────────────────────
\section{Modelo histórico de marcadores}
\label{sec:historico}

\subsection{Datos}

Utilizamos los resultados de tres ediciones del Mundial de Fútbol:
WC 2014 (Brasil), WC 2018 (Rusia) y WC 2022 (Catar), con 64 partidos
cada una, para un total de 192 partidos.
Los datos se obtienen de la base de datos de código abierto
\textit{openfootball}.

Para la \textbf{fase de grupos}, registramos el marcador al final de
los 90 minutos reglamentarios.
Para la \textbf{fase eliminatoria}, registramos el marcador al final
de los 120 minutos cuando hay tiempo extra (prórroga), dado que la
regla de puntuación de la quiniela considera todos los goles en ese
período.
Los partidos que se resuelven en tanda de penales se registran con
el marcador de empate al final del tiempo extra, pues la quiniela
instruye predecir empate en ese caso.

Mantenemos distribuciones separadas para la fase de grupos y la fase
eliminatoria, dado que los partidos eliminatorios tienden a ser más
cerrados y a producir marcadores distintos a los de la fase de grupos
\citep{groll2015}.

\subsection{Ponderación temporal: parámetro \texorpdfstring{$\gamma$}{gamma}}

Dada la evolución del nivel de juego entre ediciones del Mundial,
el historial más reciente debe pesar más en la estimación.
Definimos pesos de decaimiento geométrico: el WC 2022 tiene peso~1,
el WC 2018 tiene peso~$\gamma$, y el WC 2014 tiene peso~$\gamma^2$,
con $\gamma \in (0, 1]$.

La muestra efectiva ponderada es
\[
  n_{\mathrm{eff}} = 64 \cdot (1 + \gamma + \gamma^2).
\]

Estimamos $\gamma$ por máxima verosimilitud con validación cruzada
dejando fuera un torneo a la vez (\textit{leave-one-tournament-out}),
con un prior $\gamma \sim \mathcal{N}(1, 0.3^2)$ truncado a $(0,1]$
para regularizar hacia el caso sin decaimiento.
La estimación resultante es $\hat{\gamma} \approx 0.84$, que implica
una muestra efectiva de aproximadamente
$64 \cdot (1 + 0.84 + 0.84^2) \approx 163$ partidos ponderados.

\subsection{Suavizado mediante un prior de Poisson}

Las estimaciones empíricas de frecuencias de marcadores exactos son
inestables para muestras pequeñas: marcadores como 4-3 pueden no
haber ocurrido en 192 partidos aunque tengan probabilidad positiva.
Suavizamos la distribución histórica con un prior paramétrico.

Definimos un prior de producto de Poisson:
\[
  P_{\mathrm{prior}}(h, k) \propto \frac{\lambda_h^h}{h!} \cdot
  \frac{\lambda_k^k}{k!},
\]
donde $\lambda_h$ y $\lambda_k$ son las tasas medias de goles por
equipo estimadas a partir de los datos ponderados combinados.
La distribución suavizada combina la empírica y el prior con
$\kappa = 5$ pseudoobservaciones:
\[
  \hat{P}(h, k)
  = \frac{n_{\mathrm{eff}} \cdot P_{\mathrm{hist}}(h,k)
         + \kappa \cdot P_{\mathrm{prior}}(h,k)}
         {n_{\mathrm{eff}} + \kappa}.
\]

\subsection{Adaptación bayesiana al torneo en curso: parámetro
  \texorpdfstring{$\delta$}{delta}}

A medida que el propio WC 2026 produce resultados, esa información
debe incorporarse a la distribución.
Definimos un esquema de adaptación en el que cada partido del
WC 2026 de la jornada $j$ contribuye con peso $(1 + \delta)$ relativo
a los partidos históricos.
La distribución combinada es entonces
\[
  \hat{P}_\delta(h,k)
  \propto n_{\mathrm{eff}} \cdot \hat{P}(h,k)
        + (1 + \delta) \cdot n_{2026} \cdot P_{2026}(h,k),
\]
donde $n_{2026}$ es el número de partidos del WC 2026 ya disputados
y $P_{2026}(h,k)$ es su distribución empírica.

Estimamos $\delta$ por máximo a posteriori, dejando fuera una jornada
a la vez (\textit{leave-one-jornada-out}), con un prior
$\delta \sim \text{HalfNormal}(\sigma = 1)$ que favorece la
regularización hacia el historial cuando los datos del torneo actual
son escasos.
Al inicio del torneo, $\delta = 0$; crece conforme se acumulan
resultados y la señal del WC 2026 justifica desviarse del historial.

% ─────────────────────────────────────────────────────────────
\section{Integración de los mercados de predicción}
\label{sec:kalshi}

La distribución histórica suavizada $\hat{P}_\delta(h,k)$ captura la
estructura estadística del fútbol mundialista, pero no incorpora
información específica del partido: las fortalezas relativas de los
dos equipos, el contexto del grupo, el historial reciente.
Esa información está contenida en los precios de Kalshi.

\subsection{Corrección del sesgo favorito-longshot}
\label{sec:alpha}

Antes de usar los precios de Kalshi como probabilidades, aplicamos
una corrección por el sesgo favorito-longshot descrito en la
Sección~\ref{sec:mercados}.
Denotamos con $p_W^0$, $p_D^0$, $p_A^0$ los precios brutos de Kalshi
(normalizados para sumar 1).
La corrección sigue la transformación de potencia
\citep{whelan2024,burgi2026}:
\begin{equation}
  q_i = \frac{(p_i^0)^\alpha}{\sum_{j} (p_j^0)^\alpha},
  \qquad i \in \{W, D, A\},
  \label{eq:alpha}
\end{equation}
donde $\alpha > 1$ concentra masa en los favoritos y reduce la
sobrevaloración de los longshots.

\textbf{Inicialización.}
El valor inicial $\alpha_0 = 1.10$ está calibrado con base en el análisis
de mercados de intercambio competitivos en \cite{whelan2024}.

\textbf{Estimación secuencial.}
Tras cada jornada del Mundial, re-estimamos $\alpha$ por máximo
a posteriori~(MAP) usando todos los partidos completados
con datos de Kalshi hasta ese momento.
Para el partido $i$ con precios brutos $p_W^{0,i}, p_D^{0,i}, p_A^{0,i}$
y resultado $r_i \in \{W, D, A\}$, la log-verosimilitud de $\alpha$ es
\[
  \ell(\alpha) = \sum_{i} \log q_{r_i}(\alpha),
\]
con la prior gaussiana $\alpha \sim \mathrm{Normal}(\alpha_0,\, \sigma_\alpha^2)$,
$\sigma_\alpha = 0.20$.
El estimador MAP satisface
\[
  \hat{\alpha} = \arg\max_\alpha
    \left[\ell(\alpha) - \frac{(\alpha - \alpha_0)^2}{2\sigma_\alpha^2}\right].
\]
Cuando hay pocos datos (menos de dos partidos con resultado kalshi),
$\hat{\alpha} = \alpha_0$.
La estimación actualizada se almacena junto con $\gamma$ y $\delta$,
y se reporta en la Sección~\ref{sec:resultados}.

\subsection{Condicionamiento al resultado del partido}

Sean $p_W$, $p_D$, $p_A$ las probabilidades corregidas $q_W, q_D, q_A$
de la ecuación~\eqref{eq:alpha}.
Dividimos la distribución histórica en tres bloques:
\[
  Q_W(h,k) = \hat{P}_\delta(h,k) \cdot \mathbf{1}[h > k],
  \quad
  Q_D(h,k) = \hat{P}_\delta(h,k) \cdot \mathbf{1}[h = k],
  \quad
  Q_A(h,k) = \hat{P}_\delta(h,k) \cdot \mathbf{1}[h < k],
\]
y los reescalamos para que cada bloque sume la probabilidad de Kalshi
correspondiente:
\[
  \tilde{P}_1(h,k)
  = \frac{p_W \cdot Q_W(h,k)}{\sum_{h'>k'} Q_W(h',k')}
  + \frac{p_D \cdot Q_D(h,k)}{\sum_{h'=k'} Q_D(h',k')}
  + \frac{p_A \cdot Q_A(h,k)}{\sum_{h'<k'} Q_A(h',k')}.
\]

\subsection{Reajuste por total de goles}

Los mercados de total de goles de Kalshi proveen probabilidades
$q_n = \Prob(\text{total de goles} = n)$ para $n = 0, 1, \ldots, 6$.
Reajustamos $\tilde{P}_1$ para que las marginales de total de goles
coincidan con $q_n$:
\[
  \tilde{P}_2(h,k)
  = \tilde{P}_1(h,k) \cdot \frac{q_{h+k}}
    {\sum_{h'+k'=h+k} \tilde{P}_1(h',k')}.
\]

\subsection{Reajuste por ventaja de goles}

Análogamente, los mercados de spread de Kalshi proveen, dado que
el equipo local gana, las probabilidades $s_m^H$ de que lo haga
por exactamente $m$ goles ($m = 1, 2, 3, 4+$), y, dado que
gana el visitante, las correspondientes $s_m^A$.
Aplicamos un reajuste adicional dentro de cada bloque de resultado:
\[
  P(h,k) \propto \tilde{P}_2(h,k) \cdot r(h,k),
\]
donde $r(h,k)$ es el factor de reajuste derivado de los mercados de
spread.
La distribución resultante $P(h,k)$ es la estimación final que se
utiliza en la optimización.

% ─────────────────────────────────────────────────────────────
\section{Optimización de puntos de quiniela}
\label{sec:optimizacion}

\subsection{Regla de puntuación}

La quiniela en cuestión usa las siguientes reglas de puntuación.
Para la \textbf{fase de grupos}, con prioridad estricta
(la primera regla aplicable gana):

\begin{center}
\begin{tabular}{lc}
  \toprule
  Condición & Puntos \\
  \midrule
  Marcador exacto & 5 \\
  Ganador/empate correcto \emph{y} goles de un equipo correctos & 3 \\
  Solo ganador/empate correcto & 2 \\
  Solo goles de un equipo correctos (ganador incorrecto) & 1 \\
  Ninguna condición anterior & 0 \\
  \bottomrule
\end{tabular}
\end{center}

Para la \textbf{fase eliminatoria} (incluyendo tiempo extra; penales
equivalen a empate):
\begin{center}
\begin{tabular}{lc}
  \toprule
  Condición & Puntos \\
  \midrule
  Marcador exacto & 3 \\
  Solo ganador/empate correcto & 1 \\
  Otro & 0 \\
  \bottomrule
\end{tabular}
\end{center}

\subsection{Predicción óptima}

Dada la distribución $P(h, k)$ y la función de puntuación $f$,
los puntos esperados al predecir $(a, b)$ son
\[
  E(a, b) = \sum_{h=0}^{5} \sum_{k=0}^{5} f(a, b, h, k) \cdot P(h, k).
\]
La predicción óptima es
\[
  (\hat{a}, \hat{b})
  = \arg\max_{(a,b) \in \{0,\ldots,5\}^2} E(a, b).
\]

Nótese que $(\hat{a}, \hat{b})$ puede diferir sustancialmente del
marcador modal $\arg\max_{(h,k)} P(h,k)$.
La razón es que la regla de puntuación asigna 3 puntos a acertar
el ganador con goles de un equipo correctos, lo que favorece
marcadores del tipo $1$-$0$, $2$-$0$, $0$-$1$, $0$-$2$ sobre
marcadores bajos de empate como $0$-$0$.

\paragraph{Ejemplo: Corea del Sur vs.\ Rep.\ Checa (WC 2026,
  fase de grupos).}
Kalshi asignó probabilidades
$p_W = 0.37$, $p_D = 0.31$, $p_A = 0.32$
(local-empate-visitante), con distribución de total de goles
$\{0{:}\,9.5\%,\ 1{:}\,21\%,\ 2{:}\,29\%,\ 3{:}\,20\%,\
   4{:}\,12\%,\ 5{:}\,5\%,\ 6{:}\,3.5\%\}$.
La Tabla~\ref{tab:ejkorcze} muestra $P(a, b)$ y $E(a, b)$ para
los marcadores más probables.

\begin{table}[h]
\centering
\caption{Probabilidades $P(a,b)$ y puntos esperados $E(a,b)$
para Corea del Sur (local) vs.\ Rep.\ Checa (visita),
calculados con el modelo del artículo sin reajuste por spread.
El marcador modal y la predicción óptima difieren.}
\label{tab:ejkorcze}
\begin{tabular}{ccccc}
\toprule
Predicción $(a,b)$ & $P(a,b)$ & $E(a,b)$ & Observación \\
\midrule
$\mathbf{1\text{-}0}$ & 0.113 & $\mathbf{1.541}$ & \textbf{óptimo} \\
$2\text{-}1$ & 0.080 & 1.415 & \\
$2\text{-}0$ & 0.085 & 1.405 & \\
$1\text{-}1$ & 0.131 & 1.427 & modal (mayor $P$), empate \\
$0\text{-}1$ & 0.097 & 1.393 & \\
$0\text{-}0$ & 0.095 & 1.328 & empate \\
$1\text{-}2$ & 0.069 & 1.282 & \\
$0\text{-}2$ & 0.074 & 1.240 & \\
$2\text{-}2$ & 0.056 & 1.114 & empate \\
\bottomrule
\end{tabular}
\end{table}

El marcador modal es $1$-$1$ (empate, $P = 0.131$), pero la
predicción óptima es $1$-$0$ ($E = 1.541$).
La diferencia se debe a que $1$-$0$ también puntúa 3 puntos
cuando el resultado es $2$-$0$, $3$-$0$ o $4$-$0$, mientras que
$1$-$1$ nunca alcanza la tier de 3 puntos (véase
Observación~\ref{obs:nodraw}).
El resultado real fue $2$-$1$ (Corea del Sur).
\TODO{Este ejemplo usa únicamente el reajuste por resultado y
total de goles, sin datos de spread, para claridad expositiva.
Buscar al final del torneo un ejemplo con mayor contraste entre
el marcador modal y la predicción óptima, o donde la predicción
resultó en una puntuación especialmente alta o baja.}

\begin{observacion}[Los empates nunca son predicciones óptimas en fase de grupos]
\label{obs:nodraw}
Para cualquier predicción de empate exacto $(a, a)$, la tier de
3 puntos de la regla de puntuación de grupos es inalcanzable.
Esa tier requiere acertar simultáneamente el ganador/empate
\emph{y} los goles de al menos un equipo.
Si la predicción es $(a, a)$ y el resultado real es un empate
$(c, c)$, los goles de un equipo coinciden solo si $c = a$, caso
que corresponde al marcador exacto (5 puntos).
Para cualquier otro empate $c \neq a$, ningún gol coincide y se
obtienen 2 puntos.
Por lo tanto, los puntos esperados de predecir $(a, a)$ son
\begin{equation}
  E(a, a)
  = 5\,P(a,a)
  + 2\!\!\sum_{c \neq a} P(c,c)
  + 1 \cdot \!\!\!\!\sum_{\substack{h \neq k \\ h = a \text{ ó } k = a}}
    \!\!\!\!P(h,k),
  \label{eq:edraw}
\end{equation}
sin término de 3 puntos.
En cambio, una predicción de victoria $(a, b)$ con $a > b$ sí
accede a la tier de 3 puntos para todos los marcadores de victoria
local $(h, b)$ con $h \neq a$ ó $(a, k)$ con $k \neq b$
(resultados con un gol coincidente).
Esta asimetría estructural implica que predecir un empate puede ser
óptimo solo si la probabilidad de empate es suficientemente alta para
compensar la ausencia del término de 3 puntos.
En el fútbol mundialista, los empates ocurren en aproximadamente el
25\% de los partidos de la fase de grupos, y la probabilidad de cada
marcador de empate específico (0-0, 1-1, 2-2, \ldots) es pequeña,
de modo que este umbral no se alcanza en la práctica.
\end{observacion}

% ─────────────────────────────────────────────────────────────
\section{Implementación y plataforma web}
\label{sec:implementacion}

\subsection{Arquitectura del sistema}

El modelo se implementa en Python 3.11, organizado en un repositorio
público en GitHub.
Los componentes principales son:

\begin{itemize}
  \item \texttt{model/historical.py}: construcción de la distribución
    histórica ponderada con parámetros $\gamma$ y $\delta$.
  \item \texttt{model/kalshi.py}: obtención de precios de la API
    pública de Kalshi, corrección por sesgo favorito-longshot y
    aplicación de los tres reajustes sobre $P(h,k)$.
  \item \texttt{model/optimizer.py}: cálculo de $E(a,b)$ y argmax.
  \item \texttt{model/learn.py}: estimación de $\gamma$ y $\delta$
    por validación cruzada.
  \item \texttt{model/results.py}: obtención de resultados oficiales
    del WC 2026 desde la API pública de ESPN.
  \item \texttt{model/sanity.py}: validación automática de predicciones
    antes de su publicación (detecta predicciones degeneradas).
  \item \texttt{model/predict.py}: orquestador del pipeline completo.
\end{itemize}

La plataforma web es un sitio estático (HTML, JavaScript y CSS) sin
servidor de aplicaciones.
Las predicciones actualizadas se sirven como archivos JSON desde
GitHub Pages, lo que elimina costos de infraestructura y garantiza
alta disponibilidad.

\subsection{Pipeline de actualización automática}

Las predicciones se actualizan automáticamente mediante GitHub Actions,
el sistema de integración continua de GitHub.
El pipeline se activa en dos condiciones: por \textit{push} al
repositorio (ante cambios manuales) y por un cronograma de
\textit{cron} diseñado para cubrir todos los posibles horarios de
inicio de los 104 partidos del torneo.

Cada ejecución del pipeline sigue los pasos:
\begin{enumerate}
  \item Obtener precios de Kalshi para todos los partidos sin
    resultado (con caché de 90 minutos; los precios se congelan
    definitivamente al inicio del partido).
  \item Estimar $\gamma$ y $\delta$ actualizados.
  \item Calcular $P(h,k)$ y $(\hat{a},\hat{b})$ para cada partido.
  \item Ejecutar \texttt{model/sanity.py}; abortar si hay errores.
  \item Obtener resultados del WC 2026 desde la API de ESPN.
  \item Publicar los archivos JSON actualizados y desplegar el sitio.
\end{enumerate}

\subsection{Desafíos técnicos resueltos}

El desarrollo del sistema reveló varios problemas sutiles que
documentamos aquí, pues son representativos de los desafíos de los
sistemas de datos en tiempo real.

\textbf{Serialización JSON y claves enteras.}
Python usa enteros como claves de diccionario para las distribuciones
de goles ($\{0: 0.08, 1: 0.20, \ldots\}$), pero el formato JSON
solo admite cadenas como claves.
Al recargar el caché desde disco, las claves se recuperan como cadenas
(\texttt{"0"}, \texttt{"1"}, etc.), lo que hacía que todas las
búsquedas retornaran cero y colapsaba las predicciones al valor
degenerado 5-5.
La solución fue convertir explícitamente las claves de vuelta a enteros
al cargar el caché.

\textbf{Nombres de equipos en la fuente de resultados.}
La API de ESPN usa nombres de equipos que difieren de los nombres
oficiales FIFA en varios casos (por ejemplo, ``Czechia'' en lugar
de ``Czech Republic'', o ``Bosnia-Herzegovina'' con guion en lugar
del ampersand oficial).
El sistema de detección automática de alias alerta cuando un partido
ha superado su ventana de resultado sin que se haya importado el
marcador, permitiendo corregir rápidamente el mapeo de nombres.

\textbf{Condición de carrera en el repositorio git.}
El pipeline automático de GitHub Actions puede ejecutarse en paralelo
con cambios manuales en el repositorio, generando conflictos al
intentar hacer \textit{push}.
Se resolvió ejecutando \texttt{git pull --rebase} antes de cada
\textit{push} dentro del pipeline, lo que mantiene un historial
lineal sin conflictos de fusión.

\textbf{Entidades HTML en contextos matemáticos.}
La página del modelo matemático del sitio web usa MathJax para
renderizar expresiones LaTeX en el navegador.
El carácter ``\texttt{<}'' en expresiones como
$n_{2026,<k}$ era interpretado por el navegador como el inicio de
una etiqueta HTML antes de que MathJax pudiera procesarlo.
La solución fue usar la entidad HTML \texttt{\&lt;} dentro de las
expresiones matemáticas de la página.

\subsection{Desarrollo asistido por Claude Code}

El sistema completo se desarrolló con el apoyo de Claude Code
(anteriormente Anthropic API), una herramienta de programación
asistida por inteligencia artificial.
Claude Code actúa como un co-piloto de programación capaz de
escribir, leer, editar y ejecutar código, navegar la base de código,
y diagnosticar errores en tiempo real a partir de una descripción
en lenguaje natural del problema.

Esta modalidad de desarrollo facilitó varias tareas que
habrían requerido tiempo considerable en el ciclo tradicional:

\begin{itemize}
  \item \textbf{Arquitectura inicial}: el diseño del pipeline de
    datos (Kalshi → distribución → predicción → JSON → sitio web)
    se concretó en pocas iteraciones de especificación y refinamiento
    en lenguaje natural.
  \item \textbf{Diagnóstico de errores}: los problemas descritos en
    la sección anterior (serialización JSON, condición de carrera,
    entidades HTML) se diagnosticaron mediante la descripción del
    síntoma observable, sin necesidad de rastrear el error en el
    código manualmente.
  \item \textbf{Interfaz de usuario}: la plataforma web con soporte
    para español, banderas de países y visualización de probabilidades
    se desarrolló íntegramente a través de especificaciones en lenguaje
    natural.
  \item \textbf{Documentación}: los comentarios en el código, el
    archivo CLAUDE.md con instrucciones para el agente, y las secciones
    de metodología de la plataforma web se generaron y mantuvieron en
    paralelo con el desarrollo.
\end{itemize}

El desarrollo asistido por IA no elimina la necesidad de supervisión:
la lógica matemática del modelo, la regla de puntuación de la quiniela
y los criterios de sanidad de las predicciones requirieron revisión
cuidadosa por parte del autor.
Sin embargo, redujo el costo marginal de implementar, corregir y
documentar el código, lo que permitió dedicar más tiempo a las
decisiones de diseño estadístico.

% ─────────────────────────────────────────────────────────────
\section{Resultados}
\label{sec:resultados}

\subsection{Fase de grupos}

\TODO{Completar al concluir la fase de grupos (~2026-07-03).}

Los datos que se reportarán incluyen:

\begin{itemize}
  \item Tabla de predicciones y resultados por partido, con puntos
    obtenidos versus puntos esperados por el modelo.
  \item Puntos totales acumulados y promedio por partido.
  \item Evolución del parámetro $\delta$ a lo largo de la fase de
    grupos (número de partidos jugados vs.\ estimación de $\delta$).
  \item Valor estimado de $\hat{\alpha}$ al final de la fase de grupos
    (sesgo favorito-longshot observado en Kalshi durante el WC 2026).
  \item Comparación versus dos líneas base:
    \begin{enumerate}
      \item \emph{Línea base uniforme}: probabilidades iguales
        (1/3, 1/3, 1/3) para cada resultado.
      \item \emph{Línea base histórica}: modelo sin precios de Kalshi,
        solo distribución histórica ponderada con $\gamma$.
    \end{enumerate}
  \item Score de Brier de las probabilidades de resultado de Kalshi:
    $\text{Brier} = \frac{1}{N}\sum_{i} (p_i - o_i)^2$,
    como medida de calibración.
\end{itemize}

\subsection{Fase eliminatoria}

\TODO{Completar durante la revisión del artículo, conforme avancen
los partidos eliminatorios (~2026-07-04 al ~2026-07-19).}

Se reportarán los mismos indicadores que en la fase de grupos,
organizados por ronda: Ronda de 32, Octavos de final,
Cuartos de final, Semifinales y Final.

% ─────────────────────────────────────────────────────────────
\section{Discusión}
\label{sec:discusion}

\subsection{Aportación de los mercados de predicción}

\TODO{Completar con datos reales al concluir la fase de grupos.}

La pregunta central es si los precios de Kalshi añaden valor sobre
la distribución histórica pura.
El diseño experimental (comparación versus la línea base histórica)
permite cuantificar esta contribución.
Si el modelo con Kalshi supera consistentemente a la línea base
histórica, esto constituye evidencia de que los mercados de predicción
contienen información específica del partido que no está capturada en
el historial de Mundiales.

El modelo incluye la corrección por sesgo favorito-longshot descrita
en la Sección~\ref{sec:alpha}: los precios brutos de Kalshi se
transforman mediante la ecuación~\eqref{eq:alpha} con el exponente
$\hat{\alpha}$ estimado por MAP sobre los partidos completados con
datos de Kalshi.
Al inicio del torneo, $\hat{\alpha} = \alpha_0 = 1.10$
(valor de referencia para mercados de intercambio competitivos;
\citealp{whelan2024}).
Con el avance de la fase de grupos, $\hat{\alpha}$ se actualiza
secuencialmente para incorporar la evidencia específica de Kalshi en
el WC 2026.

\TODO{Al concluir la fase de grupos: reportar el valor estimado de
$\hat{\alpha}$ y comparar con $\alpha_0 = 1.10$. Si $\hat{\alpha}$
difiere sustancialmente, esto indica que el sesgo en Kalshi durante
el WC 2026 difiere del valor esperado para mercados de intercambio
en general, posiblemente porque el Mundial atrae a usuarios nuevos
con distinto perfil de calibración.}

\subsection{Adaptación bayesiana: el parámetro \texorpdfstring{$\delta$}{delta}}

\TODO{Analizar si $\delta > 0$ mejoró las predicciones en la fase
de grupos.}

El parámetro $\delta$ modela la hipótesis de que el WC 2026 puede
diferir sistemáticamente de los Mundiales anteriores en el número y
distribución de goles.
Si $\delta$ estimado permanece cercano a cero al final de la fase de
grupos, esto sugiere que el historial de tres Mundiales es suficiente
para caracterizar la distribución del WC 2026.
Si $\delta$ crece sustancialmente, indica una desviación real del
comportamiento histórico.

\subsection{Ausencia de predicciones de empate}

La Observación~\ref{obs:nodraw} predice que el modelo no producirá
ninguna predicción de empate en la fase de grupos.
Esta predicción se confirma empíricamente: en los 48 partidos de la
fase de grupos del WC 2026, el modelo propone en todos los casos
una victoria de uno de los dos equipos.
Este resultado no refleja una sesgo hacia las victorias en las
estimaciones de probabilidad, sino una consecuencia directa de la
regla de puntuación: la tier de 3 puntos ---que constituye una
fracción sustancial del valor esperado de una predicción--- es
estructuralmente inaccesible para los empates.
Un participante humano intuitivo que predijera empates en partidos
muy parejos estaría actuando de forma subóptima incluso si su
estimación de $P(\text{empate})$ fuera correcta.

\subsection{Limitaciones}

\textbf{Muestra histórica pequeña.}
Con solo tres Mundiales (192 partidos), la distribución histórica
tiene alta incertidumbre para marcadores infrecuentes.
El suavizado Poisson mitiga este problema pero no lo elimina.

\textbf{Truncación de la cuadrícula de marcadores.}
La optimización se realiza sobre $\{0,\ldots,5\}^2$, ignorando la
posibilidad de marcadores con 6 o más goles.
En el historial de los tres Mundiales, aproximadamente el 0.5\% de
los partidos terminaron con 6 o más goles en total.

\textbf{Variables contextuales omitidas.}
El modelo no incorpora variables específicas del partido como
lesiones, sanciones, condiciones climáticas, historial directo entre
selecciones o el contexto de clasificación dentro del grupo.
Estas variables pueden ser relevantes en algunos partidos.

\subsection{Extensiones posibles}

La arquitectura del modelo es extensible en varias direcciones.
El marco bayesiano adoptado aquí es afín a modelos que incorporan
datos de eventos dentro del partido para actualizar probabilidades
en tiempo real \citep{divekar2024}; aunque los concursos de predicción
deportiva requieren pronósticos previos al partido, esa línea de
trabajo es relevante para aplicaciones donde las predicciones pueden
ajustarse durante el juego.
El framework es aplicable a otros concursos con distintas reglas de
puntuación simplemente cambiando la función $f$.
Finalmente, el uso de mercados de predicción como fuente de prior
bayesiano es aplicable a otros torneos con menor historial que el
Mundial, donde la información histórica es más escasa.

% ─────────────────────────────────────────────────────────────
\section{Conclusiones}
\label{sec:conclusiones}

Este artículo presenta un modelo estadístico para la predicción de
marcadores exactos en el Mundial 2026 que combina dos fuentes de
información: el historial de tres ediciones del torneo (WC 2014,
2018, 2022) y los precios en tiempo real de los mercados de predicción
de Kalshi.
La integración se realiza mediante un modelo bayesiano adaptativo
con tres parámetros estimados por validación cruzada: el decaimiento
temporal $\gamma \approx 0.84$, la adaptación al torneo en curso
$\delta$, y el exponente de corrección de sesgo favorito-longshot
$\alpha$ para los precios de Kalshi.

La predicción resultante maximiza explícitamente los puntos esperados
bajo la regla de puntuación de la quiniela, que es asimétrica y
distingue varios niveles de acierto.
Esta optimización es el elemento que diferencia el modelo de un
sistema de predicción de probabilidades: no busca el marcador más
probable, sino el que rinde más en promedio bajo esa regla de puntuación
específica.

Los resultados del torneo \TODO{completar con conclusión cuantitativa
tras la fase de grupos y el análisis de calibración de Kalshi}.

El código fuente y la plataforma web están disponibles públicamente,
con actualizaciones automáticas durante todo el torneo.
Esperamos que este trabajo motive desarrollos similares para otras
competiciones y contribuya a la incipiente literatura académica en
español sobre mercados de predicción aplicados al deporte.

% ─────────────────────────────────────────────────────────────
\section*{Agradecimientos}

\TODO{Añadir agradecimientos si corresponde.}

% ─────────────────────────────────────────────────────────────
\section*{Disponibilidad del código}

El código fuente completo del modelo y la plataforma web están
disponibles en:
\begin{center}
  \url{https://github.com/alfaromurillo/quiniela}
\end{center}
bajo licencia MIT.
Las predicciones actualizadas se pueden consultar en:
\begin{center}
  \url{https://alfaromurillo.github.io/quiniela/}
\end{center}

% ─────────────────────────────────────────────────────────────
\bibliographystyle{plainnat}
\bibliography{refs}

\end{document}

%%% Local Variables:
%%% mode: LaTeX
%%% TeX-master: t
%%% End:
